\documentclass[12pt]{article}

%% Basic document formatting
\usepackage{amsmath, amsthm, amssymb, amsfonts}
\usepackage{mathtools}
\usepackage{xspace}
\usepackage{thmtools}
\usepackage{graphicx}
\usepackage{setspace}
\usepackage{fancyhdr}
\usepackage{titling}
\usepackage[left=0.4in,right=0.4in,top=1in,bottom=1in]{geometry}
\usepackage{float}
\usepackage{hyperref}
\usepackage{mathrsfs}
\usepackage{tabularx}
\usepackage[utf8]{inputenc}
\usepackage[english]{babel}
\usepackage{framed}
\usepackage[dvipsnames]{xcolor}
\usepackage{environ}
\usepackage{tcolorbox}
\tcbuselibrary{theorems,skins,breakable}

\usepackage{enumitem}
\setlist[enumerate]{leftmargin=*}
\setlist[enumerate,1]{labelindent=\parindent}
\setlist[enumerate,2]{labelindent=0pt}

% Blackboard Bold
\newcommand{\Z}{\mathbb{Z}}                 % integers
\newcommand{\Zpl}{\mathbb{Z}_{+}}           % positive integers
\newcommand{\Znn}{\mathbb{Z}_{\geq 0}}      % nonnegative integers. 
\newcommand{\Q}{\mathbb{Q}}                 % rationals
\newcommand{\Qpl}{\mathbb{Q}_{+}}           % positive Rationals
\newcommand{\R}{\mathbb{R}}                 % reals
\newcommand{\Rpl}{\mathbb{R}_{+}}           % positive Reals
\newcommand{\C}{\mathbb{C}}                 % complex numbers
\newcommand{\F}{\mathbb{F}}                 % field

% Words
\newcommand{\st}{\text{ such that }}
\newcommand{\wrt}{\text{ with respect to }}
\newcommand{\with}{\text{ with }}
\newcommand{\ie}{\text{, i.e. }}
\newcommand*{\ditto}{\text{---}\texttt{"}\text{---}}

% Operators
\newcommand{\abs}[1]{\left|#1\right|}                   % absolute value
\newcommand{\norm}[1]{\abs{\abs{#1}}}
\newcommand{\floor}[1]{\left\lfloor #1 \right\rfloor}   % floor
\newcommand{\ceil}[1]{\left\lceil #1 \right\rceil}      % ceiling
\newcommand{\powset}[1]{\mathscr{P}(#1)}

% Category Theory 
\renewcommand{\hom}{\text{Hom}}
\newcommand{\homspace}[3]{\hom_{#1}(#2, #3)}
\newcommand{\coproduct}[9]{
  \begin{tikzcd}[ampersand replacement=\&]
      \& \& #1 \arrow[dll, bend right, "#6"] \arrow[dl, "#5"] \\
   #4 \& #3 \arrow[l, "#9"] \\
      \& \& #2 \arrow[ull, bend left, "#8"] \arrow[ul, "#7"]
  \end{tikzcd}
}

\newcommand{\product}[9]{ 
  \begin{tikzcd}[ampersand replacement=\&]
      \& \& #3 \\
      #1 \arrow[r, "#7"] \arrow[urr, bend left, "#5"] \arrow[drr, bend right, "#6"] \& #2 \arrow[ur, "#8"] \arrow[dr, "#9"] \\ 
      \& \& #4
  \end{tikzcd}
}


% Algebra 
\newcommand{\Syl}[2]{\text{Syl}_{#1}{(#2)}}           % set of Sylow-p subgroups 
\newcommand{\<}{\langle}                % \<x\>, subgroup generated by x 
\renewcommand{\>}{\rangle}
\renewcommand{\index}[2]{[#1 : #2]}
\newcommand{\id}{\text{id}}             % identity element
\newcommand{\lhdneq}{%
  \mathrel{\ooalign{$\lneq$\cr\raise.22ex\hbox{$\lhd$}\cr}}}
\newcommand{\order}[1]{\text{o}(#1)}    % order of an element
\newcommand{\spec}{\operatorname{Spec}}
\newcommand{\rad}{\operatorname{rad}}   % radical of an ideal 
\newcommand{\sym}[1]{\mathfrak{S}_{#1}}
\newcommand{\aut}[1]{\text{Aut}(#1)} 
\newcommand{\gal}[2]{\text{Gal}(#1 / #2)} 
\newcommand{\inn}[1]{\text{Inn}(#1)} 
\let\oldcong\cong
\let\oldequiv\equiv
\renewcommand{\cong}{\oldequiv}
\renewcommand{\equiv}{\oldcong}

\newcommand{\triangleleftneq}{\mathrel{\ooalign{%
  $\trianglelefteq$\cr
  \hidewidth\raise0.06ex\hbox{$\not\mkern4mu$}\hidewidth\cr
}}}

\makeatletter % cycle
\newcommand{\cyc}[1]{(\mathbf{\cyc@process#1\relax})}
\def\cyc@process#1#2\relax{%
	#1%
	\ifx\relax#2\relax
	\else
		\,\cyc@process#2\relax
	\fi
}
\makeatother

\newcommand{\GL}[2]{\text{GL}_{#1}(#2)}                             % general linear group
\newcommand{\SL}[2]{\text{SL}_{#1}(#2)}                             % special linear group
\newcommand{\gl}[2]{\mathfrak{gl}_{#1}(#2)}
\renewcommand{\sl}[2]{\mathfrak{sl}_{#1}(#2)}
\newcommand{\so}[2]{\mathfrak{so}_{#1}(#2)}
\newcommand{\su}[1]{\mathfrak{su}(#1)}

% Linear Algebra
\newcommand{\bmat}[1]{\begin{bmatrix}#1\end{bmatrix}}   % bracketed matrix
\newcommand{\rank}{\operatorname{rank}}                 % rank
\newcommand{\nullity}{\operatorname{nullity}}           % nullity
\newcommand{\tr}{\operatorname{tr}}

% Combinatorics
\newcommand{\qnum}[1]{\left[#1\right]_q}                % q-analog
\newcommand{\qfact}[1]{\left[#1\right]!_q}              % q-analog factorial 
\newcommand{\qbinom}[2]{\binom{#1}{#2}_q}               % Gaussian binomial coefficient

% Topology/Analysis
\newcommand{\ball}[2]{\text{B}_{#1}(#2)}  % B_r(x): open r-balls around x
\newcommand{\diam}{\text{diam}}           % diamter of a set in metric space 
\newcommand{\lowint}[2]{\underline{\int_{#1}^{#2}}}
\newcommand{\upint}[2]{\overline{\int}_{#1}^{#2}}

% Misc Notation
\newcommand{\oneton}[1]{\{1, \ldots, #1\}}
\newcommand{\defeq}{\vcentcolon=}   % :=
\newcommand{\eqdef}{=\vcentcolon}   % =: 
\renewcommand{\bf}[1]{\textbf{#1}}
\newcommand{\im}{\operatorname{im}}
\newcommand{\fnrest}[2]{\left.#1\right|_{#2}}
\newcommand{\isom}{\overset{\sim}{\rightarrow}} 


\renewcommand{\thesection}{\arabic{section}.}
\renewcommand{\thesubsection}{(\alph{subsection})}

\newtheorem{claim}{Claim}
\newtheorem{theorem}{Theorem}
\newtheorem{proposition}{Proposition}
\newtheorem*{lemma}{Lemma}

%% Headers & title setup
\newcommand{\course}{Math140B}
\newcommand{\myname}{Jay Ser}
\setlength{\headheight}{14.5pt}
\pagestyle{fancy}
\fancyhf{}
\renewcommand{\headrulewidth}{0.4pt}
\lhead{\course}
\rhead{\myname}
\cfoot{\thepage}
\setlength{\droptitle}{-4em} 
\title{\course\ - HW \#8}
\author{\myname}
\date{2026.03.01}

\begin{document}
\maketitle
\thispagestyle{fancy}

%------------------------------------------------------------------------------%

\newcommand{\kder}[1]{\frac{d^[#1]}{dx^[#1]}}

\section*{Q1}
We have not covered the construction of the exponential function as of the due date of HW \#8, so I record the theorem in Rudin that is needed here: 
\begin{theorem}
  For any $x$,
  $$e^x = \sum_{n = 0}^{\infty} \frac{x^n}{n!}$$
  converges. 
  It follows that for every $n \in \Znn$, 
  $$\lim_{x \rightarrow \infty} x^n e^{-x} = 0$$
\end{theorem}

To solve the question, induct on $n$ for the following claim: 
for all $n \in \Zpl$, 
$$f^{(n)}(x) = \begin{cases} P_n(\frac{1}{x}) e^{-1/x^2} & \text{ if } x \neq 0 \\ 0 & \text{ if } x = 0 \end{cases}$$
where $P_n(1/x)$ is a polynomial in $1/x$. 
Consider $n = 1$. 
The $x \neq 0$ case follows by the chain rule when differentiating $e^{-1/x^2}$. 
For $x = 0$, notice that 
\begin{align*}
  f'(0) & = \lim_{x \rightarrow 0} \frac{e^{-1/x^2}}{x}  \\ 
        & = \lim_{y \rightarrow \infty} ye^{-y^2} \\ 
        & = 0
\end{align*}
where $y = 1 / x$ and the final equality follows by the same reason as Theorem 1: 
$e^{y^2} = \sum_{n = 0}^{\infty} \frac{y^{2n}}{n!} \implies y^n e^{-y^2} < \frac{(n + 1)!}{y^{n + 2}}$ for any $n$. 

Next, suppose the proposition is true for $n - 1$. 
When $x \neq 0$, $f^{(n)}$ is the product of a polynomial in $1/x$ and $e^{-1/x^2}$ by the chain rule and product rule for the derivative of $f^{(n - 1)}(x) = P_{n - 1}(1/x) e^{-1/x^2}$. 
When $x = 0$, 
\begin{align*}
  f^{(n)}(0) & = \lim_{x \rightarrow 0} \frac{P_{n - 1}(\frac{1}{x}) e^{-1/x^2}}{x} \\ 
             & = \lim_{y \rightarrow \infty} yP_{n - 1}(y) e^{-y^2} \\ 
             & = 0
\end{align*}
since each summand in the polynomial $yP_{n - 1}(y)$ multiplied by $e^{-y^2}$ goes to 0 as $y \rightarrow \infty$. 
This completes the induction and thereby completes the problem. 

\section*{Q6} 
\subsection{} 
Fix any $y \in \R$. Differentiate with respect to $x$ to obtain $f(x + y) = f(x) f(y)$
$$f'(x + y) = f(y) f'(x)$$ 
Hence, letting $x = 0$, 
\begin{equation} \label{eq:idk}
  f'(y) = f'(0) f(y)
\end{equation}
for any $y \in \R$. 
Let $c = f'(0)$ and consider $g(x) = \frac{f(x)}{e^{cx}}$. 
Differentiating, 
$$g'(x) = \frac{e^{cx} f'(x) - ce^{cx} f(x)}{e^{2cx}} = \frac{f'(x) - cf(x)}{e^{ct}}$$ 
But \eqref{eq:idk} says the denominator of $g'(x)$ is zero for all $x \in \R$, hence 
$$g'(x) = 0 \text{ for all } x \in \R$$
It follows that $f(x) / e^{cx}$ is a constant, so $f(x) = k e^{cx}$.
Finally, $f(0 + y) = f(0)f(y)$ and the fact that $f(y) \neq 0$ for some $y \in \R$ means $f(0) = k e^0 = k = 1$, so 
$$f(x) = e^{cx}$$ 

\subsection{} 
Let $g(x) = \log(f(x))$. 
Then 
\begin{equation} \label{eq:hom} 
  g(x + y) = \log(f(x)f(y)) = \log(f(x)) + \log(f(y)) = g(x) + g(y)
\end{equation}
i.e., $g$ is an endomorphism on $\R$. 
Note $g$ is indeed defined on all of $\R$ because $f > 0$ by the following: 
$f(x) = f(x / 2 + x / 2) = (f(x / 2))^2$, so $f(x) \geq 0$ on all of $\R$. 
But also $f(0) = f(x) f(-x)$ and $f(0) = 1$ as shown in (a), hence $f(x) \neq 0$ on all of $\R$.  
Because $g$ is an endomorphism, for any $m / n \in \Q$, 
\begin{equation} \label{eq:cnt} 
  g(m) = n g(\frac{m}{n}) \implies g(\frac{m}{n}) = \frac{1}{n} g(m) = \frac{1}{n} m g(1) = \frac{m}{n} g(1)
\end{equation}
In other words, denoting $g(1) = c$, $g(q) = cq$ for all $q \in Q$. 
Finally, because $f$ and $\log$ are continuous, $g$ is continuous, and $\Q$ dense in $\R$, 
$$g(x) = cx \text{ for all } x \in \R$$ 
By the relationship between $f$ and $g$, it follows that $f(x) = e^{cx}$ on all of $\R$. 


\end{document}
